\textbf{Phần 1. Trắc nghiệm lựa chọn}
\setcounterpageref{ex}{0}

\begin{ex} %[3D1N05]
Kết quả của phép tính $3 \times 6$ là
\choice
{$16$}
{$21$}
{$24$}
{\True $18$}
\loigiai{
$3 \times 6 = 18$.\\
\textbf{Đáp án của câu hỏi là D}
}
\end{ex}

\begin{ex} %[3D1N05]
Kết quả của phép tính $21 : 3$ là
\choice
{\True $7$}
{$6$}
{$8$}
{$9$}
\loigiai{
$21 : 3 = 7$.\\
\textbf{Đáp án của câu hỏi là A}
}
\end{ex}

\begin{ex} %[3D1H05]
Kết quả của phép tính $3 \times 9$ là
\choice
{$24$}
{$30$}
{\True $27$}
{$21$}
\loigiai{
$3 \times 9 = 27$.\\
\textbf{Đáp án của câu hỏi là C}
}
\end{ex}

\begin{ex} %[3D1H05]
Kết quả của phép chia $27 : 3$ là
\choice
{\True $9$}
{$8$}
{$7$}
{$6$}
\loigiai{
$27 : 3 = 9$.\\
\textbf{Đáp án của câu hỏi là A}
}
\end{ex}

\begin{ex} %[3D1H05]
Điền số thích hợp vào chỗ chấm: $3 \times \dots = 24$.
\choice
{$7$}
{\True $8$}
{$9$}
{$6$}
\loigiai{
$24 : 3 = 8$.\\
\textbf{Đáp án của câu hỏi là B}
}
\end{ex}

\begin{ex} %[3D1H05]
Điền số thích hợp vào chỗ chấm: $\dots : 3 = 5$.
\choice
{$12$}
{$18$}
{\True $15$}
{$20$}
\loigiai{
$5 \times 3 = 15$.\\
\textbf{Đáp án của câu hỏi là C}
}
\end{ex}

\begin{ex} %[3D1V05]
Tính giá trị biểu thức: $3 \times 8 + 16$.
\choice
{$38$}
{\True $40$}
{$42$}
{$45$}
\loigiai{
$3 \times 8 + 16 = 24 + 16 = 40$.\\
\textbf{Đáp án của câu hỏi là B}
}
\end{ex}

\begin{ex} %[3D1V05]
Tính giá trị biểu thức: $30 : 3 + 45$.
\choice
{\True $55$}
{$50$}
{$60$}
{$65$}
\loigiai{
$30 : 3 + 45 = 10 + 45 = 55$.\\
\textbf{Đáp án của câu hỏi là A}
}
\end{ex}

\begin{ex} %[3D1V05]
Mỗi nhóm học tập có 3 bạn. Hỏi 7 nhóm như thế có bao nhiêu bạn?
\choice
{$18$ bạn}
{\True $21$ bạn}
{$24$ bạn}
{$27$ bạn}
\loigiai{
$3 \times 7 = 21$ bạn.\\
\textbf{Đáp án của câu hỏi là B}
}
\end{ex}

\begin{ex} %[3D1V05]
Có 24 cái kẹo chia đều cho 3 bạn. Mỗi bạn nhận được số cái kẹo là
\choice
{$7$ cái}
{$9$ cái}
{$6$ cái}
{\True $8$ cái}
\loigiai{
$24 : 3 = 8$ cái kẹo.\\
\textbf{Đáp án của câu hỏi là D}
}
\end{ex}

\begin{ex} %[3D1V05]
Mẹ mua 3 hộp sữa chua, mỗi hộp có 4 hộp nhỏ. Tổng số hộp sữa chua nhỏ mẹ mua là
\choice
{$10$ hộp}
{$14$ hộp}
{\True $12$ hộp}
{$16$ hộp}
\loigiai{
$4 \times 3 = 12$ hộp.\\
\textbf{Đáp án của câu hỏi là C}
}
\end{ex}

\begin{ex} %[3D1C05]
Tìm một số biết rằng lấy số đó nhân với 3 rồi trừ đi 5 thì được 19.
\choice
{$7$}
{$9$}
{$6$}
{\True $8$}
\loigiai{
Số đó nhân 3 bằng: $19 + 5 = 24$.\\
Số đó là: $24 : 3 = 8$.\\
\textbf{Đáp án của câu hỏi là D}
}
\end{ex}

\textbf{Phần 2. Trắc nghiệm đúng sai}
\setcounterpageref{ex}{0}

\begin{ex} %[3D1H05]
Xét tính đúng sai của các phép tính trong bảng nhân 3, bảng chia 3:
\choiceTF
{\True $3 \times 4 = 12$.}
{\True $18 : 3 = 6$.}
{$3 \times 7 = 24$.}
{\True $3 \times 7 = 21$.}
\loigiai{
\begin{itemize}
    \item $3 \times 4 = 12$: Đúng.
    \item $18 : 3 = 6$: Đúng.
    \item $3 \times 7 = 21 \neq 24$. Ý 3 Sai, Ý 4 Đúng.
\end{itemize}
\begin{itemize}
    \item a) \textbf{Đúng}.
    \item b) \textbf{Đúng}.
    \item c) \textbf{Sai}.
    \item d) \textbf{Đúng}.
\end{itemize}
}
\end{ex}

\begin{ex} %[3D1V05]
Bác An có 18 quả cam chia đều vào 3 túi. Xét tính đúng sai:
\choiceTF
{\True Mỗi túi có 6 quả cam.}
{\True Nếu ăn hết 1 túi thì còn lại 12 quả cam.}
{Nếu chia vào mỗi túi 9 quả cam thì được 3 túi.}
{\True Nếu chia vào mỗi túi 9 quả cam thì được 2 túi.}
\loigiai{
\begin{itemize}
    \item $18 : 3 = 6$ quả: Đúng.
    \item Còn lại $6 \times 2 = 12$ quả: Đúng.
    \item $18 : 9 = 2$ túi $\neq 3$ túi. Ý 3 Sai, Ý 4 Đúng.
\end{itemize}
\begin{itemize}
    \item a) \textbf{Đúng}.
    \item b) \textbf{Đúng}.
    \item c) \textbf{Sai}.
    \item d) \textbf{Đúng}.
\end{itemize}
}
\end{ex}

\textbf{Phần 3. Trả lời ngắn}
\setcounterpageref{ex}{0}

\begin{ex} %[3D1H05]
Tính tích: $3 \times 5$.
\shortans{15}
\loigiai{
\textbf{Đáp số:} $15$\\
$3 \times 5 = 15$.
}
\end{ex}

\begin{ex} %[3D1H05]
Tính thương: $12 : 3$.
\shortans{4}
\loigiai{
\textbf{Đáp số:} $4$\\
$12 : 3 = 4$.
}
\end{ex}

\begin{ex} %[3D1V05]
Tính giá trị biểu thức: $3 \times 7 + 9$.
\shortans{30}
\loigiai{
\textbf{Đáp số:} $30$\\
$21 + 9 = 30$.
}
\end{ex}

\begin{ex} %[3D1V05]
Có 27 bông hoa cắm đều vào 3 bình. Hỏi mỗi bình có bao nhiêu bông hoa?
\shortans{9}
\loigiai{
\textbf{Đáp số:} $9$\\
$27 : 3 = 9$ bông hoa.
}
\end{ex}

\begin{ex} %[3D1V05]
Một tam giác đều có 3 cạnh bằng nhau, mỗi cạnh dài $8\text{ cm}$. Chu vi tam giác đó là bao nhiêu xăng-ti-mét?
\shortans{24}
\loigiai{
\textbf{Đáp số:} $24$\\
$8 \times 3 = 24\text{ cm}$.
}
\end{ex}

\begin{ex} %[3D1C05]
Tìm $x$ biết: $x \times 3 = 15 + 12$.
\shortans{9}
\loigiai{
\textbf{Đáp số:} $9$\\
$x \times 3 = 27 \Rightarrow x = 27 : 3 = 9$.
}
\end{ex}
